\chapter{Poker Without Revealing the Cards}

\section{Objective}
This document accompanies the executable poker demonstrator and presents the
same construction in a more formal order. The implementation does not yet
attempt to solve the full game of Texas Hold'em; it isolates the dealing
problem into four stages that can be inspected in code and in runtime logs.

\section{Scientific Scope}
The present demonstrator addresses the following narrow question:

\begin{quote}
Can mutually distrustful parties instantiate a poker hand without a trusted
dealer while preserving privacy of hole cards until showdown?
\end{quote}

The current answer is partial rather than complete. The demonstrator models:
\begin{enumerate}
    \item randomization of a canonical deck,
    \item secure multi-party shuffling,
    \item secure dealing of private and public cards,
    \item public commitment to private hands.
\end{enumerate}

It does not yet model betting, hand-ranking evaluation, or adversarial network
execution.

\section{Source Mapping}
\begin{itemize}
    \item C source directory:
    \path{/src/src/poker_without_revealing_the_cards/}
    \item Primary implementation file:
    \texttt{poker\_without\_revealing\_the\_cards.c}
    \item Build entrypoint: \texttt{demostrations/Makefile}
    \item PDF helper: \texttt{./book -pdf section\_poker}
\end{itemize}

\section{Protocol Decomposition}
\subsection*{Stage 1: Deck Randomization}
A standard 52-card deck is permuted and a public digest of the resulting state
is published. The digest fixes the transcript without revealing the deck order.

\subsection*{Stage 2: Secure Multi-Party Shuffling}
Each player privately shuffles the deck and applies a commutative-encryption
layer. After each transformation, the encrypted deck is committed to publicly.
Once the final encrypted deck is fixed, the demonstrator also publishes a hash
for each encrypted card that will later be dealt. These public hashes are the
visible form of the sampled cards during this stage.

\subsection*{Stage 3: Secure Dealing}
Hole cards are privately recovered by removing encryption layers in a
player-specific order. Community cards are jointly decrypted and become public.
Before each decryption, the demonstrator verifies that the drawn encrypted card
matches the Stage~2 hash commitment for that deal slot. It also validates that
no plaintext card appears twice.

\subsection*{Stage 4: Hand Commitments}
If player \(P_i\) holds hole cards \(h_{i,1}, h_{i,2}\), the player samples a
nonce \(N_i\) and publishes
\[
    C_i = H\!\left(N_i \parallel 2 \parallel h_{i,1} \parallel h_{i,2}\right).
\]
Later, \(N_i\) and the private cards are revealed so that every participant can
recompute \(C_i\).

\section{Build and Run}
\begin{verbatim}
cd /src
make -C demostrations demo_poker
\end{verbatim}

Alternative suit rendering modes:
\begin{verbatim}
make -C demostrations demo_poker POKER_ARGS="--suits ascii"
make -C demostrations demo_poker POKER_ARGS="--suits unicode"
POKER_SUIT_MODE=unicode make -C demostrations demo_poker
\end{verbatim}

\section{Observed Runtime Structure}
The program emits a staged trace of the form:
\begin{verbatim}
Stage 1  Randomization of the deck
Stage 2  Secure multi-party shuffling
Stage 3  Secure multi-party dealing
Stage 4  Hand commitments and validation
\end{verbatim}
The log also records public commitments, encrypted deck previews, dealt cards,
sampled-card hash commitments, sampled-card verification, and commitment
verification at showdown.

\section{Limitations}
The current demonstrator should be interpreted as a protocol sketch with an
executable transcript, not as a final cryptographic poker engine. Open work
includes:
\begin{enumerate}
    \item betting rounds and action synchronization,
    \item formal hand evaluation,
    \item authenticated transport,
    \item malicious-security guarantees,
    \item stronger proof obligations for correct shuffle and deal operations.
\end{enumerate}
